\section{The Business Starts Before the Business}

\subsection{The Pre-Existing Inventory}

When people describe early founders as starting from nothing, they usually mean that there is no company revenue, no office, no investor, and no product. But that is not the same as having no resources. Founders may own human capital (skills and judgment acquired through learning and work) alongside social capital, credibility, savings, language ability, cultural knowledge, access, and available attention.

Saras Sarasvathy's work on \textit{effectuation} offers a useful starting point: begin with who you are, what you know, and whom you know, rather than assuming that you must first predict a distant market and collect all the resources necessary to serve it. Effectuation also emphasizes affordable loss, stakeholder commitments, and adaptation to unexpected events.\footnote{Saras D. Sarasvathy, ``Causation and Effectuation: Toward a Theoretical Shift from Economic Inevitability to Entrepreneurial Contingency,'' \textit{Academy of Management Review} 26, no. 2 (2001): 243--263.} It is a valuable lens, not a guarantee of success or a replacement for market validation. Researchers have actively debated its explanatory boundaries.\footnote{John~Kitching and Julia~Rouse, ``Contesting effectuation theory: Why it does not explain new venture creation,'' \textit{International Small Business Journal} 38, no.~6 (2020): 515--536.}

The founder's opening inventory is larger than the traditional 3 questions. I would examine at least 10 kinds of starting assets:

\begin{tcolorbox}[
  enhanced,
  colback=white,
  colframe=BayesBorder!80!white,
  boxrule=0.4pt,
  arc=2pt,
  left=3pt,right=3pt,top=3pt,bottom=3pt,
  before skip=3pt,after skip=5pt
]
\scriptsize
\setlength{\tabcolsep}{2.5pt}
\renewcommand{\arraystretch}{1.15}
\begin{tabularx}{\linewidth}{@{} >{\bfseries\color{BayesInk}}p{1.9cm} >{\color{BayesInk}}X >{\color{BayesSecondary}}X @{}}
\toprule
\textbf{Starting Asset} & \textbf{Operational Meaning} & \textbf{The Core Test} \\
\midrule
Skills & Capabilities you reliably execute & Can you perform to a commercial standard? \\
Knowledge & Models of how a domain functions & Does insight reveal non-obvious gaps? \\
Credibility & Valid reasons buyers trust claims & Will a buyer risk money on your promise? \\
Relationships & People willing to exchange help & Will they respond and commit resources? \\
Distribution & Routes to reach an audience & Can you contact buyers with permission? \\
Capital & Cash available for bounded tests & What can you risk without harm to family? \\
Time and focus & Predictable weekly working blocks & Can you sustain deep work for months? \\
Tools, systems & Software, IP, hardware, workflows & Do these cut delivery time or cost? \\
Constraints & Commitments, debt, legal gates & What paths are closed or dangerous? \\
Motivation & Problems you persist in solving & Will you work after novelty fades? \\
\bottomrule
\end{tabularx}
\end{tcolorbox}

The word \textit{asset} does not mean a company owns these things, or that every one has a market price. A relationship is not your property. Professional experience does not entitle you to an employer's confidential information. A professional contact list does not automatically give you permission to advertise. These distinctions are essential, because what you can recall, what you can use, and what you can legally commercialize are different categories.

The inventory is also not static. Today you may have no reputation among hospital administrators. After 20 serious conversations and 1 useful pilot, you may have substantially better insight and more credible access. \textbf{The starting asset is not a sentence about your identity. It is the current state of a system capable of learning and accumulating new assets.}

\subsection{The 3 Layers of Each Asset}

Every asset should be examined through 3 separate questions:

\begin{enumerate}[label=\textbf{\arabic*.},leftmargin=11pt,itemsep=1.5pt,topsep=1.5pt,parsep=0pt]
  \item \textbf{Possession:} Is the capability, knowledge, relationship, or resource actually present?
  \item \textbf{Activation:} Can you deploy it now, in this context, under real constraints?
  \item \textbf{Conversion:} Does deploying it create an advantage a buyer, partner, or employee values?
\end{enumerate}

A founder may know a university dean but have no relationship strong enough to earn a meeting. That is possible possession without proven activation. They may earn a meeting but discover that the dean cannot purchase anything. That is activation without an obvious route to commercial conversion.

Similarly, an engineer may build sophisticated models but face enterprise buyers whose primary complaint is sluggish customer support and who prefer a basic workflow tool. Technical strength is real; commercial relevance is unproved.

A useful shorthand is:
\begin{center}
\textbf{Potential asset $\rightarrow$ usable asset $\rightarrow$ economically consequential asset.}
\end{center}

The arrows require evidence. Many founder mistakes happen when the first state is mistaken for the third.
